\chapter{Perception for Intelligent Surveillance}
\label{chap:background_perception}

\section{Introduction}

Every surveillance decision ultimately rests on perception: converting raw pixels and audio samples into structured, semantically meaningful events. This chapter reviews the four perception capabilities used by AURA-MAS edge agents---object detection, multi-object tracking, video anomaly detection, and audio event detection---with emphasis on the accuracy/latency trade-offs that govern edge deployment.

\section{Object Detection}

Object detection localizes and classifies objects in images. Modern real-time detection is dominated by the \gls{yolo} family, which frames detection as a single dense regression pass over the image. The current generation (YOLO11, released by Ultralytics in 2024 \cite{ultralytics2025yolov8}) provides a spectrum of model scales; the nano variant (YOLO11n, $\approx$2.6M parameters) achieves approximately 39.5\% \gls{map} on COCO while sustaining real-time throughput on CPU and embedded hardware, making it the natural choice for per-camera edge agents. Transformer-based detectors such as RT-DETR \cite{ultralytics2025rtdetr} eliminate non-maximum suppression and reach higher accuracy at moderate additional cost; they constitute a drop-in upgrade path in our architecture since agents encapsulate their detector behind a common event interface.

Two practical considerations matter for surveillance specifically. First, detection confidence thresholds trade recall against precision at the \emph{event} level, not just the box level: a missed person is a missed intrusion. Second, pretrained COCO classes (person, car, truck, backpack, \ldots) cover the vast majority of surveillance-relevant objects, so fine-tuning is optional rather than blocking---an important de-risking property for a time-boxed project \cite{encord2024yolo}.

For zone-based rules, a bounding box is not itself a location model. AURA-MAS
uses the bottom centre of a detected box (the foot point) and a ray-casting
test against the declared polygon. This is a practical image-plane heuristic:
it is often more stable than the box centre for a standing person, but it can
fail under occlusion, a steep top-down view, a cropped body, or a detector box
that includes substantial background. The IoU used for static-object logic is
defined by
\[
\operatorname{IoU}(A,B)=\frac{|A\cap B|}{|A\cup B|}.
\]
An IoU threshold of 0.6 has different operational meaning for a small backpack
than for a truck-sized box; scale, detector jitter, and camera perspective
change the same numerical overlap. This is why the thresholds in
Appendix~\ref{app:config} are disclosed as implementation parameters rather
than universal surveillance constants.

\section{Multi-Object Tracking}

Detection alone yields per-frame boxes; surveillance rules (loitering, line crossing, abandonment) require \emph{identity over time}. Multi-object tracking (MOT) associates detections across frames into trajectories. The prevailing paradigm is tracking-by-detection, and the reference algorithm is ByteTrack \cite{zhang2022bytetrack}, whose key insight is to associate \emph{every} detection box---including low-confidence ones---in a two-stage matching (high-confidence first, then low-confidence against unmatched tracks), which markedly reduces identity switches under occlusion. ByteTrack is integrated natively in the Ultralytics runtime (\texttt{model.track(..., tracker="bytetrack.yaml")}), which AURA-MAS uses directly. Evaluation of trackers uses HOTA/MOTA/IDF1 metrics computed by tools such as TrackEval \cite{luiten2020trackeval, mendez2024metrics}; multi-camera extensions of tracking (cross-camera re-identification) exist \cite{amosa2023multicamera, fei2023momct} but are deliberately excluded from this thesis on privacy grounds---our cross-camera corroboration operates at the \emph{event} level, not the identity level.

ByteTrack's second association stage is especially relevant at the deployed
5-FPS sampling rate: low-confidence detections can bridge a short occlusion,
but sparse temporal sampling gives an object more time to move between frames
and therefore makes association harder. The corpus contains no per-frame track
annotations, so HOTA, MOTA, and IDF1 are named here as appropriate tracker
metrics but are deliberately not reported as results. A field-of-view
homography would provide a more principled relation between camera views and
ground-plane zones than image-plane polygons; it is also the missing input
needed to make the auction's overlap term meaningful.

\section{Video Anomaly Detection}

\gls{vad} aims to flag frames or segments that deviate from normal patterns. Three families coexist \cite{gong2025survey}:

\begin{itemize}
    \item \textbf{Reconstruction/prediction methods} (autoencoders, future-frame prediction) trained on normal data only; effective on constrained scenes (Avenue, ShanghaiTech) but brittle under domain shift.
    \item \textbf{Weakly-supervised methods}, trained with video-level labels on datasets such as UCF-Crime \cite{sultani2018ucfcrime}; the current accuracy leaders.
    \item \textbf{Zero-shot semantic methods}, which exploit vision-language alignment: a frame embedding is compared against natural-language descriptions of normal versus anomalous situations using \gls{clip} \cite{radford2021clip}. Recent variants (VadCLIP-style, AVadCLIP with audio \cite{wu2025avadclip}, training-free EventVAD \cite{shao2025eventvad}, explainable EX-VAD \cite{huang2025exvad}) show that useful anomaly signals can be obtained \emph{without any training}, which is exceptionally attractive for a one-month project.
\end{itemize}

AURA-MAS uses the third family as an optional semantic scorer on sampled frames: anomaly score is the softmax mass assigned to a set of anomalous text prompts (``people fighting'', ``fire and smoke'', \ldots) relative to normal prompts. This complements---rather than replaces---the deterministic zone-rule engine, and every \gls{clip}-triggered event is subject to the same downstream fusion, verification, and policy thresholds as any other event.

\section{Audio Event Detection}

Audio is a cheap, omnidirectional, lighting-invariant modality that video-only systems ignore. Sound event detection (\gls{sed}) classifies acoustic events---glass breaking, screams, gunshots, alarms---from short audio windows. Pretrained models dominate practice: YAMNet (MobileNet-based, 521 AudioSet classes) \cite{tensorflow_yamnet, hershey2017cnn}, PANNs \cite{kong2019panns}, the Audio Spectrogram Transformer \cite{li2023ast}, and BEATs \cite{chen2022beats} offer increasing accuracy at increasing cost; efficient distillations exist for edge deployment \cite{schmid2024improving, saradopoulos2025real}. A surveillance-relevant subset of classes can be mapped directly to alert-worthy event types, with per-class confidence thresholds. Because deep audio stacks can be heavy for constrained hardware, AURA-MAS additionally implements a dependency-light \gls{dsp} fallback: rolling z-scores of short-time energy and spectral flatness flag acoustic transients as generic \texttt{audio\_anomaly} events---sufficient to demonstrate multimodal fusion even where TensorFlow is unavailable.

The audio repair documented for the v2 campaign illustrates why temporal
pre-processing is part of the method, not incidental plumbing. Audio is
represented as a log-mel time--frequency signal for model inference. A
one-second chunk is evaluated with a lookback extension and YAMNet scores are
max-pooled across the model's internal frames; using mean pooling on
independently re-windowed one-second chunks diluted brief transients below
threshold. The before/after observation and the campaign implications are
recorded in Appendix~\ref{app:methodology}.

\section{Multimodal Fusion}

Fusion strategies are classified as early (feature-level), late (decision-level), or hybrid. For a distributed system, \emph{late fusion} is the natural fit: each agent emits calibrated event confidences, and a fusion node combines them without needing access to raw features---preserving both bandwidth and privacy. AURA-MAS groups events into spatio-temporal hypotheses (same incident family, same zone, within a sliding window) and combines confidences with a noisy-OR model weighted by per-modality reliability, plus small calibrated bonuses for cross-modality and cross-sensor corroboration (Section~\ref{sec:design_fusion}). Audio-visual \gls{vad} research confirms that such corroboration materially improves robustness \cite{wu2025avadclip, ustun2026multimodal}.

Noisy-OR is interpretable, but its usual probabilistic reading assumes that
the contributing evidence is conditionally independent. Repeated events from
one camera watching the same subject violate that assumption. In the present
implementation they can still drive the product toward one, so monotonicity
should not be conflated with calibration. The fusion audit in
Chapter~\ref{chap:evaluation} treats this saturation as a limitation rather
than as a multimodal gain.

\section{Datasets and Evaluation Data}

Public datasets relevant to this work include UCF-Crime (1{,}900 untrimmed real-world videos across 13 anomaly classes) \cite{sultani2018ucfcrime}, Avenue and ShanghaiTech for scene-specific \gls{vad}, VIRAT for activity detection from static surveillance cameras, and MEVA for large-scale multi-camera activity detection \cite{tan2025how}. Synthetic data generation is an increasingly credible complement \cite{delussu2024synthetic, riaz2023synthetic, chu2025lfsyngen}. For \emph{system-level} evaluation (as opposed to model-level), no standard benchmark exists that combines multi-sensor streams with ground-truth incident annotations; this thesis therefore contributes replayable scenario manifests built from public clips and scripted recordings (Chapter~\ref{chap:evaluation}), following the methodological lead of ETISEO \cite{nghiem2007etiseo} and modern MLOps evaluation practice \cite{aicity2025}.

\section{Conclusion}

Pretrained detectors, trackers, zero-shot anomaly scorers, and audio classifiers can be combined without training a model for this project. The remaining challenge is how to fuse, verify, decide on, and explain their outputs. The rest of the thesis describes that architecture.
